\documentclass[a4paper]{article}
\usepackage{amsfonts}
\usepackage{amsmath}
\usepackage{amssymb}
\usepackage{graphicx}     

\begin{document}
  
\noindent \large Auteur: Victoria Mazur \\
\noindent \large Studierichting: Informatica\\
\noindent \large Oefeningenreeks 1E nr. 34\\

\medskip

\normalsize

Gegeven: $A(z) = (z - 1)(z + 1)(z - z_3)$ waar $z_3$ onbekend is. \\

We weten dat $A(i) = 4$. \\

Substitutie van $z = i$ in de vergelijking: \\

$A(i) = (i - 1)(i + 1)(i - z_3) = 4$ \\

Eerst berekenen we $(i - 1)(i + 1)$: \\

$(i - 1)(i + 1) = i^2 - 1^2$ \\

Omdat $i^2 = -1$: \\

$(i - 1)(i + 1) = -1 - 1 = -2$ \\

Nu substitueren we dit terug: \\

$-2(i - z_3) = 4$ \\

Beide zijden delen door $-2$: \\

$i - z_3 = \dfrac{4}{-2} = -2$ \\

Oplossen naar $z_3$: \\

$z_3 = i - (-2) = i + 2 = 2 + i$ \\

Verificatie door $A(0)$ te berekenen: \\

$A(0) = (0 - 1)(0 + 1)(0 - z_3)$ \\

$A(0) = (-1)(1)(0 - (2 + i))$ \\

$A(0) = (-1)(1)(-(2 + i))$ \\

$A(0) = (-1)(-(2 + i))$ \\

$A(0) = 2 + i$ \\

Dus de complete vorm van $A(z)$ is: \\

$A(z) = (z - 1)(z + 1)(z - 2 - i)$ \\

\begin{thebibliography}{999}
%\bibitem{a} Referentie
\end{thebibliography}
\end{document}